\section{Arithmetic towers and equivariant gluing at a node}
\label{sec:parameters}

We construct the parameter to which the unramified input of
Proposition~\ref{prop:equivariant-unramified} will be applied.  There are
two steps.  First we lift the arithmetic parameter $\rho$ to a punctured
curve in characteristic zero, retaining its Frobenius action and its
entire boundary homomorphism.  We then place this lift on one component
of a separating nodal curve and extend it to an unramified parameter on
a smooth generic fiber.  In the second step the gluing must preserve
both Frobenius and compactness of the monodromy image.  A ramified
auxiliary component will allow us to satisfy these two conditions
simultaneously.

\subsection{Lifting the arithmetic parameter}

With $k=\overline{\F}_q$, set
\begin{equation}\label{eq:coefficient-traits}
 R_0=W(\F_q),\qquad D=\Frac(R_0),\qquad
 R=W(k),\qquad K=\overline{\Frac(R)},\qquad
 \Dbar=\overline D\subset K.
\end{equation}
Here $\Dbar$ is the algebraic closure of $D$ inside $K$.  The ring $R$
is an excellent complete strictly henselian discrete valuation ring
with residue field $k$.  Extend the Witt-vector automorphism of $R$
induced by $a\mapsto a^q$ on $k$ to an automorphism $\varphi$ of $K$.
It fixes $D$ and preserves $\Dbar$.  On roots of unity of $\ell$-power
order it acts by $\zeta\mapsto\zeta^q$, as follows from their
Teichm\"uller lifts in $R$.  We use the same letter for these compatible
actions, with the coordinate orientation fixed in
Section~\ref{sec:conventions}.

Choose a smooth projective pointed lift
\[
 (\mathscr X,\mathfrak x)\longrightarrow\Spec R_0
 \quad\text{of }(X_0,x_0),
 \qquad \mathscr U=\mathscr X\setminus\mathfrak x.
\]
Pointed deformations are unobstructed because
$H^2(X_0,T_{X_0}(-x_0))=0$; an ample line bundle also lifts, and formal
algebraization gives the projective family.  This is the lifting
construction of \cite[Section~8.5]{CT}, here performed over $R_0$ so as
to retain the field of definition.  Write
\[
 (C_1,x_1)=(\mathscr X,\mathfrak x)_D,
 \qquad U_1=C_1\setminus\{x_1\}.
\]
By smoothness, a function $t$ on a neighborhood of $\mathfrak x$ over
$R_0$ may be chosen as a relative parameter.  Its differential
trivializes the conormal line of the section.  We also regard $t$ as a
rational function on $C_1$.

For an integer $d$ invertible in $R_0$, the coordinate $t$ identifies
the residual gerbe of $\mathscr X(\sqrt[d]{\mathfrak x})$ with
$B\mu_d$ over $R_0$: it trivializes the divisor line by $1/t$ and
thereby specifies its trivial $d$th root.  Pulling a finite torsor
back to this neutral object gives its boundary fiber.  A frame over
the strictly henselian base $R$ identifies that fiber with its finite
structure group, acting on itself on the right.  In this frame, both inertia
and semilinear transport are recorded by left multiplication.  The
construction below chooses such frames compatibly through the tower.

\begin{proposition}\label{prop:lifted-parameter}
The arithmetic representation $\rho$ determines a continuous
parameter
\[
 \sigma_1:\pi_1^{\mathrm{et}}((U_1)_{\Dbar})
             \longrightarrow\PGL_n(L')
\]
for a finite coefficient extension $L'/L$, together with equivariance
under $\varphi$, with the following properties.
\begin{enumerate}[(i)]
\item Its geometric image is Zariski dense.  After compatible choices
of frames it is the same compact subgroup as
$\rho(\pi_1^{\mathrm{et}}(U_{0,k}))$.
\item Its boundary homomorphism is
\begin{equation}\label{eq:lifted-boundary}
 \gamma\longmapsto\exp(t_\ell(\gamma)N_1),
\end{equation}
where $N_1$ is regular nilpotent.  Every prime factor of
characteristic-zero inertia other than $\ell$ acts trivially.
\item In compatible frames at the coordinate-neutralized root gerbes
of $x_1$, the equivariance has a single value
$A\in H_0$, where $H_0=\rho(\pi_1^{\mathrm{et}}(U_0))$, and
\begin{equation}\label{eq:frobenius-scales-nilpotent}
                    \Ad(A)N_1=qN_1.
\end{equation}
\item For every algebraic representation of $\PGL_n$, the associated
local system over $(U_1)_K$ admits an invariant lattice whose finite
reductions extend lisse over $\mathscr U_R$, equivariantly under
$\varphi$.  Their specializations recover the tensor local system of
$\rho$, including its given arithmetic action.
\end{enumerate}
\end{proposition}

We prove the proposition after a finite-cover lemma that will also be
used in the nodal smoothing.  It compares monodromy orbits on finite
torsor fibers, allowing disconnected torsors as required when we
enlarge the compact group that records monodromy.

\begin{lemma}\label{lem:component-orbits}
Let $S$ be the spectrum of a complete strictly henselian discrete
valuation ring with algebraically closed residue field, and let
$\mathcal Y\to S$ be a proper flat tame Deligne--Mumford curve with
finite diagonal, projective coarse space, and geometrically reduced
fibers.  Let $\mathcal T\to\mathcal Y$ be a finite \'etale right
$Q$-torsor, for a finite group $Q$, and let $y:S\to\mathcal Y$ be a
section in its smooth scheme locus.  Identify the geometric special
and generic fibers of $y^*\mathcal T$ by a trivialization over $S$.
The partitions of this finite set by connected components of
$\mathcal T_{\bar s}$ and $\mathcal T_{\bar\eta}$ then agree.
\end{lemma}

\begin{proof}
The stack $\mathcal T$ is again proper, flat, and tame, with
geometrically reduced fibers.  Proper henselian invariance lifts the
open-and-closed decomposition of $\mathcal T_{\bar s}$ to one of
$\mathcal T$; equivalently, it lifts the corresponding idempotents
\cite[Theorem~4.5]{Rydh}.  Consider one lifted summand $\mathcal Z$.
Its special fiber is connected, reduced, and proper over an
algebraically closed field, so
\[
             H^0(\mathcal Z_{\bar s},\mathcal O)=k(\bar s).
\]
Coarse-space pushforward for a tame stack is exact and commutes with
base change \cite[Lemmas~2.3.3--2.3.4]{AbramovichVistoli}.  The coarse space
here is flat over $S$: locally, taking
invariants under a linearly reductive finite group preserves flatness.
Thus semicontinuity on this proper coarse space gives
$h^0(\mathcal Z_{\bar\eta},\mathcal O)\leq1$.  Flatness ensures that
the generic fiber is nonempty, so this dimension is exactly one.
In particular $\mathcal Z_{\bar\eta}$ is connected.  Every lifted
special component therefore has exactly one geometric generic
component.  Intersecting these components with the trivialized
section fiber proves the assertion.
\end{proof}

\begin{proof}[Proof of Proposition~\ref{prop:lifted-parameter}]
The use of the arithmetic image $H_0$, rather than only the geometric
image, retains the action that we need to specialize at the end of the
proof.  This image is compact.  Choose a faithful linear
representation and an $H_0$-invariant lattice, and let
$H_{0,r}$ be a decreasing sequence of open normal congruence subgroups
cofinal at the identity.  Put $\Gamma_r=H_0/H_{0,r}$.  The finite
quotients of $\rho$ give a compatible tower of right $\Gamma_r$-torsors
on $U_0$.  We choose the torsor convention so that associated local
systems have monodromy $\rho$.

At stage $r$ the inertia image has $\ell$-power order.  Choose indices
$d_r$, each a power of $\ell$, with $d_r\mid d_{r+1}$ and with $d_r$
divisible by that order.  The stage-$r$ torsor extends uniquely to a
finite \'etale torsor on the root stack
\[
                 X_0(\sqrt[d_r]{x_0}).
\]
Indeed a Kummer root of the local parameter kills the finite tame
inertia, after which the cover is \'etale across the root divisor.
This local description also proves the assertion over the original
finite field, or one can descend the geometric extension by its
uniqueness.  The tame-cover interpretation is
\cite[Proposition~A.10]{LieblichOlsson}.

The relative root stack
$\mathscr X(\sqrt[d_r]{\mathfrak x})$ is smooth, proper, and tame over
$R_0$, with finite diagonal, since $d_r$ is invertible in $R_0$.
Proper henselian invariance therefore lifts the torsor to this
relative root stack \cite[Theorem~4.5]{Rydh}; see also
\cite[Theorem~A.11 and Corollaries~A.12--A.13]{LieblichOlsson} and
\cite[Sections~8.3 and~8.5]{CT}.  Full faithfulness lifts its group
action and the transition maps, comparing on a divisible root index
when necessary, and preserves all relations among those maps.  Only
the root indices have been required to be prime to $p$; the finite
groups $\Gamma_r$ may have order divisible by $p$.

Restriction to $\mathscr U$ and then to $(U_1)_{\Dbar}$ gives the
parameter $\sigma_1$ with values in
$\varprojlim_r\Gamma_r=H_0$, after framing.  All the lifted data are
over $R_0$.  Their base changes to $R$ and $\Dbar$ consequently carry
the descended semilinear $\varphi$-actions and their compatible
transition maps.

We first check the geometric image.  Choose a section of
$\mathscr U_R$ through a geometric special point and compatible
frames of the tower along it.  Such frames exist because this section
is strictly henselian.  Apply Lemma~\ref{lem:component-orbits} at
each stage after base change to $R$.  The geometric fibers of the
root stacks and of their finite \'etale covers are smooth curves as
stacks.  A connected component is therefore irreducible, and removing
the marked gerbes and their inverse images preserves its
connectedness.  The component partitions on the punctured geometric
special and generic fibers agree.  In the fiber of a torsor these
partitions are the monodromy orbits; the orbit of the chosen frame
is precisely the image subgroup of $\Gamma_r$.  Thus the geometric
images have the same image in every $\Gamma_r$.  Both are compact,
hence closed, in $H_0$, so they are equal.  Connected components are
unchanged by extension from $\Dbar$ to $K$, giving the assertion for
$\sigma_1$ itself.  Density now follows from the hypothesis on $\rho$.

To compare the boundary actions, use $t$ to neutralize the residual
gerbe of each root stack as $B\mu_{d_r}$, compatibly under power
maps.  Pullback to its neutral object over $R$ gives a finite
\'etale torsor on a strictly henselian trait.  The resulting finite
nonempty sets of frames have surjective transition maps, so a compatible
frame can be chosen throughout the tower.  Its $\mu_{d_r}$-action is
the same on special and generic fibers under the identification of
prime-to-$p$ roots of unity.  These actions record the entire
peripheral homomorphism, not just its conjugacy class.  Passing to
the limit gives \eqref{eq:lifted-boundary}, with $N_1$ conjugate to
the original regular nilpotent after the fixed generator convention.
Since only $\ell$-power root indices were used, each
$\Z_a(1)$-factor of characteristic-zero inertia with $a\ne\ell$
acts trivially, including the factor $a=p$.

The neutral objects defined by $t$ are themselves defined over $R_0$.
In their compatible frames the semilinear action is left
multiplication by elements of $\Gamma_r$ compatible in $r$, hence
by one element $A\in H_0$.  Its compatibility with the gerbe action
reads
\[
 A\,\exp(t_\ell(\gamma)N_1)\,A^{-1}
      =\exp(q\,t_\ell(\gamma)N_1).
\]
Taking the logarithm of a nontrivial boundary element proves
\eqref{eq:frobenius-scales-nilpotent}.  All these statements use the
chosen coordinate orientation.  Replacing the Frobenius generator
$\varphi$ by $\varphi^{-1}$ replaces $A$ by $A^{-1}$ and $q$ by
$q^{-1}$.

Finally, an algebraic representation defined over a finite coefficient
extension has a lattice invariant under the compact group $H_0$.
Each finite reduction factors through some $\Gamma_r$, so the lifted
torsor extends that reduction over $\mathscr U_R$.  These are
equivariant extensions, and their special fibers are the reductions
of the original arithmetic local system.  Tensor products and
morphisms compare through the same tower; for morphisms one may
clear denominators in the chosen lattices.  This gives the tensor
specialization data in (iv).
\end{proof}

\subsection{A compact conjugator and an auxiliary component}

The first parameter is now lifted with its arithmetic action.  To
extend it across a separating node, the two boundary actions must
agree when expressed using the opposite orientations of the two
branches.  The following linear-algebra observation gives an
identification that also commutes with $A$ and belongs to a fixed
compact group.

\begin{lemma}\label{lem:compact-conjugator}
Suppose $N_1\in\mathfrak{pgl}_n(L')$ is regular nilpotent,
$A\in\PGL_n(L')$ satisfies \eqref{eq:frobenius-scales-nilpotent}, and
$H_0\subset\PGL_n(L')$ is compact with $A\in H_0$.  After a finite
coefficient extension there exist a cocharacter
$c:\mathbb G_m\to\PGL_n$, a compact open subgroup $J$ containing
$H_0$, and an integer $e>1$ prime to $\ell$, such that
\begin{equation}\label{eq:compact-conjugator}
 c(z)A=Ac(z),\qquad \Ad(c(z))N_1=zN_1,
 \qquad P:=c(-1/e)\in J.
\end{equation}
Consequently $P$ commutes with $A$ and
\begin{equation}\label{eq:inverse-boundary-conjugation}
 P\exp(-e aN_1)P^{-1}=\exp(aN_1)
 \qquad(a\in L').
\end{equation}
\end{lemma}

\begin{proof}
Identify $N_1$ with its traceless nilpotent matrix and choose a matrix
lift $\widetilde A$ of $A$.  The relation
$\widetilde A N_1=qN_1\widetilde A$ preserves the regular Jordan
filtration.  On its successive one-dimensional quotients the
eigenvalues of $\widetilde A$ form a geometric progression with ratio
$q$.  They are distinct, since $q>1$ is an integer and the coefficient
field has characteristic zero.  Thus $\widetilde A$ is diagonalizable
after a finite extension.  Choose an eigenvector $v_0$ generating a
Jordan chain, and put $v_i=N_1^iv_0$ for $0\leq i<n$.  For some
$a_0\ne0$ these vectors satisfy
\[
 \widetilde A v_i=a_0q^i v_i,
 \qquad N_1v_i=v_{i+1}\ (i<n-1),\qquad N_1v_{n-1}=0.
\]
The projective class of the diagonal cocharacter
$c(z)v_i=z^iv_i$ commutes with $A$ and scales $N_1$ by $z$.

Choose a closed faithful linear representation of $\PGL_n$ and a
lattice $\Lambda$ stable under $H_0$.  The inverse image of
$\GL(\Lambda)$ is a compact open subgroup $J$ of $\PGL_n(L')$
containing $H_0$: it is open by continuity and compact because the
embedding is closed.  Since $c(1)=1$ and $J$ contains a neighborhood
of $1$, choose $e>1$ with $e\equiv-1$ modulo a sufficiently high
power of $\ell$.  Then $-1/e$ is sufficiently close to $1$ that
$c(-1/e)\in J$.  This choice makes $e$ prime to $\ell$.
The last identity follows by applying
$\Ad(P)N_1=-N_1/e$ to the exponential.
\end{proof}

Fix $c,J,e,P$ as in the lemma, enlarging the coefficient field once
and again denoting it by $L'$.  Retain the notation for all parameters.
Define a cover
\begin{equation}\label{eq:auxiliary-cover}
 h:C_2\longrightarrow C_1
 \quad\text{by normalizing }C_1\text{ in }D(C_1)(w),
 \qquad w^e=t.
\end{equation}
The valuation of $t$ at $x_1$ is one.  The polynomial $w^e-t$ is
therefore Eisenstein there, even after extension to $\Dbar$, so this
is a degree-$e$ geometrically connected cover.  The curve $C_2$ is
smooth and projective over $D$, and there is a unique point $x_2$
over $x_1$.  It is $D$-rational and totally ramified, with local
parameter $w$; these assertions are also visible in the completed
local extension $D[[t]]\subset D[[w]]$.  Riemann--Hurwitz gives
$g(C_2)\geq2$.

Put $U_2=C_2\setminus\{x_2\}$ and
$\sigma_2=h^*\sigma_1$ on $(U_2)_{\Dbar}$.  Any other branch points
of $h$ lie over $U_1$ and cause no ramification in this pullback
local system.  In the coordinate $w$, its boundary homomorphism is
$\exp(e\,t_\ell(\gamma)N_1)$.

For the rest of the construction choose open normal subgroups
$J_r\subset J$, decreasing and cofinal at $1$, and set
\begin{equation}\label{eq:nodal-finite-quotients}
                         Q_r=J/J_r.
\end{equation}
The map $H_0\to Q_r$ factors through one of the earlier quotients
$\Gamma_j$, because $H_0\cap J_r$ is open in $H_0$.  Extension of
structure group therefore gives a right $Q_r$-torsor on the first
component with its original equivariance; pullback by $h$ gives the
one on the second.  These torsors need not be connected.  For any common
root index $b$ divisible by their stage-$r$ inertia orders they extend
over $C_i(\sqrt[b]{x_i})_{\Dbar}$.  The cover $h$ induces a map of
these root stacks, since $h^*x_1=e x_2$.  In coordinates it sends a
$b$th root of $t$ to the $e$th power of a $b$th root of $w$; on the
residual gerbes its map is
\begin{equation}\label{eq:gerbe-power-map}
                  \mu_b\longrightarrow\mu_b,
                  \qquad \zeta\longmapsto\zeta^e.
\end{equation}
The neutral objects fixed by $t$ and $w$ correspond under this map:
in their divisor-line trivializations,
$h^*(1/t)=(1/w)^e$, with no scalar factor.  The second torsor and its
frame are pulled back from the first, and this fiber identification is
defined over $D$.  Thus the induced frame on the second gerbe has exactly
the same equivariance value $A$ as the first.  This equality, in addition to
the factor $e$ in the inertia, is the arithmetic information supplied
by the algebraic cover.

\subsection{The balanced smoothing and its parameter}

Identify $x_1$ with $x_2$ and write $C_0=C_1\cup_{x_1=x_2}C_2$.
Choose a projective smoothing
$\mathscr C\to\Spec D[[s]]$ of this separating node, with a
relatively ample line bundle.  The curve has unobstructed deformations
and an independent smoothing parameter at the split node; lifting an
ample line bundle algebraizes the deformation.  As in
\cite[Section~8.2]{CT}, the completed local equation can be chosen as
\begin{equation}\label{eq:ordinary-node}
                           a b'=s,
\end{equation}
where $a$ restricts to a parameter on the first branch and $b'$ to
one on the second.  Rescaling their tangent coordinates over $D$, and
rescaling $s$ accordingly, ensures
\[
             (a/t)(x_1)=1,\qquad (b'/w)(x_2)=1.
\]
These equalities will identify the coordinate neutralizations of the
branch gerbes without a constant unit twist.

Choose compatible roots $s_b$ of $s$ and put
\begin{equation}\label{eq:nodal-traits}
 S_b=\Spec\Dbar[[s_b]],\qquad s_b^b=s,
 \qquad \Omega=\overline{\Dbar((s))}.
\end{equation}
Extend $\varphi|_{\Dbar}$ coefficientwise to $\Dbar((s))$, then to
the union obtained by adjoining the compatible roots $s_b$ while
fixing those roots, and finally to $\Omega$.  We again denote this
automorphism by $\varphi$.  The smooth projective geometric generic fiber
of $\mathscr C$ will be denoted $C/\Omega$.  It is connected of genus
$g(C_1)+g(C_2)$, and its polarization is $\varphi$-invariant.

For each $b$, let $\mathcal C(b)\to S_b$ be the balanced twisted
model with node chart
\begin{equation}\label{eq:balanced-node}
 \left[\Spec\bigl(\Dbar[[s_b]][u,v]/(uv-s_b)\bigr)/\mu_b\right],
 \quad
 \begin{cases}
 a=u^b,\quad b'=v^b,\\
 \zeta(u,v)=(\zeta u,\zeta^{-1}v).
 \end{cases}
\end{equation}
Away from the closed node this is the ordinary family after base
change.  The construction is the balanced twisted-curve construction
of \cite[Theorems~1.9--1.10, Remark~1.11, and
Proposition~2.2(ii)]{OlssonTwisted}; the chart and its power maps are
described in \cite[Section~8.2]{CT}.  In particular each model is
proper, flat, and tame, with projective coarse space and geometrically
reduced fibers.  Its special-fiber normalization is the disjoint
union of the two root stacks $C_i(\sqrt[b]{x_i})_{\Dbar}$.
For $b\mid m$ the power maps on $u,v$ give the corresponding maps
between models over the map $S_m\to S_b$.

The construction may be made over $D[[s_b]]$ with the group scheme
$\mu_b$, so the models and these maps carry the chosen semilinear
actions.  The ratios $a/t$ and $b'/w$ have residue one.  Their formal
roots with constant term one are unique and compatible under power
maps in characteristic zero.  They are also preserved by $\varphi$.
Consequently the node chart identifies the branch neutral objects
with those already used for $t,w$.  For general $b$, only the
$\ell$-primary quotient of the gerbe acts on our torsors.  On that
quotient $\varphi$ acts by the $q$th power; the other prime factors
act trivially on the torsors.

We fix at this point the index that will produce Borel level on the
special-fiber bundle stack.  Choose a strictly dominant cocharacter
$\lambda\in X_*(T)$ for the split torus $T\subset B\subset G$ and an
integer $b_*$ satisfying
\begin{equation}\label{eq:alcove-index}
       0<\langle\alpha,\lambda\rangle<b_*
       \qquad\text{for every positive root }\alpha.
\end{equation}
Choose the parameter indices $b_r$ so that
$b_*\mid b_r\mid b_{r+1}$ and $b_r$ kills the two stage-$r$ inertia
actions.  The sheaf construction will use the fixed model
$\mathcal C(b_*)$ and inertia type $\lambda|_{\mu_{b_*}}$; the varying
indices $b_r$ supply only finite-stage extensions of the parameter.  Denote by
$\mathcal V\subset\mathcal C(b_*)$ its smooth scheme locus.  Its
geometric generic fiber is $C$ and its special fiber is
$(U_1\sqcup U_2)_{\Dbar}$.

\begin{proposition}\label{prop:nodal-parameter}
There exists a continuous Zariski-dense parameter
\[
              \tau:\pi_1^{\mathrm{et}}(C)\longrightarrow J
                      \subset\PGL_n(L')
\]
with equivariance under $\varphi$.  For every algebraic representation
of $\PGL_n$, an invariant lattice in its evaluation local system has
finite reductions extending lisse over
$\mathcal V\times_{S_{b_*}}S_{b_r}$ for sufficiently large $r$.
These extensions are equivariant and specialize on the two punctured
components to $\sigma_1$ and $\sigma_2$, with their given actions,
compatibly with the entire tensor system.
\end{proposition}

\begin{proof}
We first glue the finite torsors on the special fiber of
$\mathcal C(b_r)$.  Express both actions using the common node
generator $\zeta\in\mu_{b_r}$.  Write $u_{1,r}(\zeta)\in Q_r$ for
the first action in its fixed frame.  On the second branch the
positive coordinate generator is $\zeta^{-1}$.  Combining this
orientation reversal with \eqref{eq:gerbe-power-map}, its action in
the induced frame is $u_{1,r}(\zeta)^{-e}$.  If $P_r$ is the image
of $P$ in $Q_r$, Equation~\eqref{eq:inverse-boundary-conjugation}
gives
\begin{equation}\label{eq:finite-stage-gluing}
       P_r\,u_{1,r}(\zeta)^{-e}P_r^{-1}=u_{1,r}(\zeta).
\end{equation}
To read this identity directly at a finite root index, lift
$\zeta$ to characteristic-zero tame inertia and use
\eqref{eq:lifted-boundary} before reducing to $Q_r$.  Different
lifts have the same reduction.  The prime factors other than $\ell$
contribute the identity on both sides.

Identify the second right-torsor fiber with the first by left
multiplication by $P_r$.  Equation~\eqref{eq:finite-stage-gluing}
says exactly that this identification intertwines their gerbe
actions.  Both semilinear actions have value $A_r$, the reduction of
$A$.  Since $P_rA_r=A_rP_r$, the same identification intertwines
Frobenius.  In coordinates on the two fibers the required square
is the elementary equality
\[
                P_r(A_rg)=A_r(P_rg)\qquad(g\in Q_r).
\]
Figure~\ref{fig:node-gluing} records the two signs and this gluing
map.  All the identifications are reductions of the one element
$P\in J$, so they are compatible throughout the tower.

\begin{figure}[!b]
\centering
\begin{tikzpicture}[>=Stealth,thick,font=\small]
  \coordinate (n) at (0,0);
  \draw[linkblue] (-4.6,0.2) .. controls (-2.7,0.8) and (-1.3,0.5) .. (n);
  \draw[linkblue] (n) .. controls (1.3,-0.5) and (2.7,-0.8) .. (4.6,-0.2);
  \fill (n) circle (2pt);
  \node[above=5pt] at (-3.2,0.45) {$(C_1,x_1)$};
  \node[below=5pt] at (3.2,-0.45) {$(C_2,x_2)$};
  \node[above=7pt] at (n) {$B\mu_{b_r}$};
  \node[align=center] at (-2.8,-0.75)
    {$u\mapsto\zeta u$\\$u_{1,r}(\zeta)$};
  \node[align=center] at (2.8,0.75)
    {$v\mapsto\zeta^{-1}v$\\$u_{1,r}(\zeta)^{-e}$};
  \node[draw,rounded corners,inner sep=6pt] (f1) at (-2.8,-2.05)
    {first fiber: $Q_r$};
  \node[draw,rounded corners,inner sep=6pt] (f2) at (2.8,-2.05)
    {second fiber: $Q_r$};
  \draw[->] (f2.west) -- node[above] {$g\mapsto P_rg$}
    node[below] {$P_rA_r=A_rP_r$} (f1.east);
\end{tikzpicture}
\caption{The two branches of the balanced node and the torsor fibers on
its normalization, shown schematically.  The cover $w^e=t$ multiplies
boundary monodromy by $e$, and the second branch reverses its orientation.  Multiplication
by $P_r$ identifies the resulting action $u_{1,r}^{-e}$ with
$u_{1,r}$ while commuting with the common Frobenius value $A_r$.}
\label{fig:node-gluing}
\end{figure}


The normalization torsors now glue to a finite \'etale right
$Q_r$-torsor $\mathcal T_{r,0}$ on $\mathcal C(b_r)_0$.  This is
ordinary nodal gluing in the equivariant chart: \'etale locally on
each branch cover, the torsor is constant across the origin; its
two constant fibers have been identified with the same
$\mu_{b_r}$-action.  Glue them across $uv=0$ and take the
equivariant quotient.  This proves \'etaleness at the node as well
as away from it, and retains the $Q_r$-action and semilinear action.
It is the finite-cover gluing of \cite[Section~8.3]{CT}, with the
identification now chosen to commute with $A$.

Proper henselian invariance lifts $\mathcal T_{r,0}$ to a finite
\'etale torsor $\mathcal T_r$ on $\mathcal C(b_r)$
\cite[Theorem~4.5]{Rydh}.  Full faithfulness lifts the group actions,
the equivariance isomorphisms, and the transition maps after pullback
to divisible indices.  Their relations lift as well.  Restriction
to the common geometric generic curve $C$ yields a tower of finite
$Q_r$-torsors and hence, after compatible framing, a continuous
parameter with image in
$\varprojlim_r Q_r=J$.  The lifted equivariance gives the asserted
action on this parameter.

We verify density using the first component, without imposing
connectedness on these $Q_r$-torsors.  Choose a section of
$\mathcal V\to S_{b_*}$ through a point of $(U_1)_{\Dbar}$,
base-change it to every $S_{b_r}$, and choose compatible frames of the
lifted tower along these sections.  Apply
Lemma~\ref{lem:component-orbits} over each $S_{b_r}$.  The orbit of
the frame under the monodromy of the punctured first component lies
in its connected component in the whole special torsor.  The lemma
identifies that component's section fiber with the geometric generic
component's section fiber, which is the generic monodromy orbit.
Thus, if $H_1$ is the image of $\sigma_1$ and $H_\tau$ the image of
$\tau$ in these compatible frames, then
\[
        \operatorname{im}(H_1\to Q_r)
          \subset\operatorname{im}(H_\tau\to Q_r)
        \qquad\text{for every }r.
\]
Both subgroups are compact and therefore closed in $J$.  Since the
$J_r$ are cofinal, these inclusions imply $H_1\subset H_\tau$.
Proposition~\ref{prop:lifted-parameter} makes $H_1$ Zariski dense,
and hence $\tau$ is Zariski dense.

Finally let $V$ be an algebraic representation over a finite
coefficient extension and choose a $J$-stable lattice in $V$.
Every finite reduction factors through some $Q_r$.  Away from the
node the models $\mathcal C(b_r)$ agree with the corresponding
base changes of $\mathcal C(b_*)$.  The torsor $\mathcal T_r$
therefore supplies the required lisse extension on
$\mathcal V\times_{S_{b_*}}S_{b_r}$, and its special restrictions
are exactly the two component torsors with their semilinear actions.
Using this one tower for all representations supplies the tensor and
morphism comparisons, as in Proposition~\ref{prop:lifted-parameter}.
These are the equivariant lattice-specialization data needed for
nearby cycles on the fixed model $\mathcal C(b_*)$.
\end{proof}
